\chapter{Related Work}

\section{3D Gaussian Splatting Foundations} 
3D Gaussian Splatting (3DGS) shifts the rendering paradigm from the volumetric ray-marching common in NeRFs to a rasterization-based approach. The scene is represented by a collection of 3D Gaussians, where each primitive is defined by its position, covariance, opacity, and spherical harmonic coefficients. By projecting these primitives onto the image plane and sorting them for alpha-blending, 3DGS achieves real-time rendering speeds while preserving high-frequency details~\cite{kerbl20233dgs}. Standard 3DGS initializes these primitives using sparse point clouds derived from Structure-from-Motion (SfM). During optimization, the method employs an adaptive density control strategy: gradients are monitored to split Gaussians in under-reconstructed areas and prune those with low opacity. However, because the optimization is driven primarily by photometric loss, the resulting geometry often suffers from artifacts. Without explicit geometric regularization, the underlying surface representation tends to be noisy or fragmented, even if the rendered novel views appear visually pleasing.

\section{Initialization Strategies} 
The fidelity of the final reconstruction is heavily dependent on the quality of the initial point cloud. Standard 3DGS creates a strict dependency on sparse SfM data (typically from COLMAP)~\cite{schoenberger2016sfm}, which can act as a bottleneck for challenging captures or large-scale scenes. To mitigate this, methods like GeoGaussian move away from sparse seeds, instead utilizing dense depth and normal priors to align Gaussian orientations with the physical surface from the onset~\cite{li2024geogaussian}. This geometry-aware initialization improves convergence and surface consistency compared to photometry alone.

Recent work by Roman Mutel offers a distinct alternative by bypassing the SfM pipeline entirely. This approach processes RGBD streams from consumer devices (such as the iPhone Pro) and utilizes Visual-Inertial Odometry (VIO) for pose estimation, effectively unprojecting depth maps directly into 3D space. By doing so, the method achieves a linear time complexity \(O(n)\), a significant improvement over the typical polynomial complexity associated with global bundle adjustment, and demonstrates that VIO poses can be sufficiently accurate for training even when traditional feature matching fails.

\section{Depth and Normal Regularization} 
Initialization alone cannot prevent geometric drift during the optimization process. In regions with weak texture or repetitive patterns, photometric loss is an insufficient constraint, leading to “floating” artifacts. DN-Splatter addresses this ill-posed problem by integrating strictly enforced geometric priors. Rather than relying on mesh extraction, it utilizes monocular depth and normal estimates to supervise the training. The key contribution of such methods is the adaptability of the regularization: stronger constraints are applied in areas of high geometric uncertainty, while photometric evidence takes precedence where image data is clear~\cite{dn_splatter}. This balance significantly reduces surface noise, though it primarily focuses on local consistency rather than enforcing a global topology.

\section{Surface-Based Constraints and Structural Representations}
To ensure global geometric coherence, recent methods have attempted to bind Gaussians to explicit surface representations. SuGaR (Surface-Aligned Gaussian Splatting) constrains the primitives to align with a mesh extracted from the density field. By enforcing this alignment, the method acts as a bridge between volumetric rendering and traditional mesh-based graphics, facilitating physics simulations and cleaner surface exports~\cite{sugar}. MiLO extends this idea by placing the mesh “in the loop,” allowing the mesh topology to actively guide the splitting and pruning of Gaussians dynamically~\cite{milo2024}.
A closely related direction explicitly targets planar and structured surfaces, which are especially prominent in indoor scenes. PlanarGS~\cite{planargs} introduces planar primitive modeling into the 3DGS framework by detecting dominant planar surfaces and jointly optimizing both planar parameters and Gaussian primitives. This planar prior enforces stronger global constraints on large flat surfaces such as walls, floors, and ceilings, improving flatness and normal consistency while reducing redundant Gaussians in large planar regions. However, this approach depends on robust plane detection and is primarily tailored to structured indoor geometry.
Another closely related line of work focuses explicitly on surface reconstruction quality from Gaussian representations. PGSR (Point-Guided Gaussian Surface Reconstruction) proposes a Gaussian-based pipeline tailored for high-fidelity surface extraction rather than purely novel view synthesis. The method reconstructs surfaces by deriving depth maps from Gaussian splatting via opacity-weighted integration and subsequently fusing them into a volumetric representation using TSDF. By explicitly optimizing for surface accuracy and leveraging point-guided constraints, PGSR achieves improved surface completeness and normal consistency compared to standard 3DGS pipelines. However, similar to other surface-oriented approaches, PGSR relies on post-hoc depth integration and volumetric fusion, rather than enforcing surface-level constraints directly during Gaussian optimization, which can limit its ability to maintain global structural consistency in large planar indoor scenes~\cite{pgsr}.
IndoorGS further explores structure-aware Gaussian splatting specifically for indoor environments by introducing layout-aware priors and indoor-specific geometric regularization. The method exploits common indoor assumptions such as dominant orthogonal planes and consistent room layouts to stabilize geometry and reduce surface fragmentation. By incorporating scene-level structural cues, IndoorGS demonstrates improved reconstruction consistency in large indoor spaces compared to generic 3DGS pipelines~\cite{indoorgs}.
Another important recent work is GausSurf, which integrates multi-view geometry guidance into the Gaussian splatting optimization process to enhance surface reconstruction quality. GausSurf distinguishes between texture-rich and texture-less regions, leveraging PatchMatch-based multi-view stereo (MVS) geometry for the former and normal priors for the latter. This geometry-guided strategy improves both surface detail and optimization stability, leading to higher-quality surface reconstruction and more efficient training compared to purely photometric or single-modality geometric approaches~\cite{gcaussurf}.
Alternatively, SAGS (Structure-Aware Gaussian Splatting) avoids explicit meshing in favor of implicit structural encoding. It employs a graph-based representation to learn relationships between neighboring primitives, enforcing meaningful spatial coherence without the overhead of managing an explicit mesh. These approaches highlight a trade-off in current literature: mesh-based methods offer tighter geometric coherence and explicit surface extraction, while implicit structural representations provide more compact geometry with reduced memory footprints.

\section{Research Gap} 
Current literature largely treats geometric priors in isolation. Methods tend to focus exclusively on one modality—either improved initialization, depth and normal regularization, planar or structural constraints—making it difficult to assess how these strategies interact. It remains unclear, for example, whether dense initialization renders depth regularization redundant, or if surface constraints can function effectively without high-quality structural priors such as those provided by multi-view geometry integration.

To address this, we propose a unified framework that disentangles these components. By integrating initialization strategies inspired by GeoGaussian~\cite{li2024geogaussian}, adaptive regularization from DN-Splatter~\cite{dn_splatter}, planar surface constraints as in PlanarGS~\cite{planargs}, and mesh-based alignment similar to SuGaR~\cite{sugar} into a single modular pipeline, we allow for the independent activation of each modality. This approach enables a systematic ablation study to understand the interplay and specific trade-offs between normals, depth, and structural constraints in 3D Gaussian Splatting, particularly for challenging indoor planar surfaces.
